\chapter{Produit scalaire, orthogonalité et projections}
\section*{Exercice 6.1 -- Angle}
Le produit scalaire vaut $1$, les deux normes valent $\sqrt2$. Ainsi $\cos\theta=1/2$ et $\theta=\pi/3$. Le résultat est un angle non orienté entre $0$ et $\pi$.

\section*{Exercice 6.2 -- Projection sur une droite}
Le coefficient de projection est $2/2=1$, donc $H=(1,1,0)$. $MH^2=1+1+1=3$, d'où distance $\sqrt3$. Le produit $\overrightarrow{HM}\cdot(1,1,0)=0$ vérifie le calcul.

\section*{Exercice 6.3 -- Normalité}
Les deux produits scalaires valent zéro. Les deux directions ne sont pas colinéaires ; toute combinaison est donc orthogonale à $(1,1,1)$ par bilinéarité. Cela prouve la normalité au plan entier.

\section*{Exercice 6.4 -- Deux plans}
Le produit scalaire des normales vaut $2-2+0=0$. Les normales étant non nulles, les plans sont perpendiculaires, quelles que soient leurs constantes dans les équations cartésiennes.

\section*{Exercice 6.5 -- Minimum}
$MH$ est orthogonal au plan, donc à $HK$. Pythagore donne $MK^2=MH^2+HK^2>MH^2$ puisque $HK>0$. Les distances sont positives ou nulles ; prendre les racines donne la stricte inégalité.

\section*{Exercice 6.6 -- Inégalité triangulaire}
Développer le carré : $\|u+v\|^2=\|u\|^2+2u\cdot v+\|v\|^2\le(\|u\|+\|v\|)^2$. Les deux membres comparés sont positifs ou nuls ; la racine carrée conserve l'ordre.
